\documentclass[12pt]{article}

% 使用 ctex 包来支持中文
\usepackage{ctex}

% 用于数学公式
\usepackage{amsmath}
\usepackage{amssymb}

% 页面边距设置
\usepackage[a4paper, margin=1in]{geometry}

% 定义标题、作者和日期
\title{关于物质局域化、基本粒子和其相互作用的思考}
\author{} % 作者留空
%\date{}   % 日期留空

\begin{document}

\maketitle

\begin{abstract}

在本文中，我们将继续在《计算实在论》这一我们正在构建的理论框架下，对``物质''这一最基本的物理实在，作一次深入而开放的展开思考。
我们希望读者能将这篇文章，视为一份坦诚的``研究日志''。我们的目的，是邀请您与我们一同踏上这段逻辑推理的旅程。在思考过程中，我们将会发现，一些看似简单的原理，可能会以一种近乎必然的方式，推导出一个个有趣且深刻的现象与结论。

我们无意宣称这里提出的模型是最终的答案。我们深知，任何一个试图从根基上重构整个粒子物理学的叙事，都必然会面临最严苛的审视，甚至可能被归类于非主流的探索。我们坦然接受这一切。
我们分享这些思考的唯一目的，是希望这个统一且逻辑自洽的视角，能够为您——我们尊敬的同行探索者——提供一些新的、或许有益的思考角度。如果文中的任何一个推演环节，能为您带来片刻的、关于现实本质的智识上的共鸣，那么我们工作的价值，便已得到了最好的肯定。

\end{abstract}

\section{局域化原理——从混沌中涌现}

在《计算实在论》的框架下，所有基本的存在——无论是作为``结构''的物质（孤子），还是作为``过程''的光（光子）——其能够作为一个稳定的、局域化的实体而存在，都源于同一个深刻的动力学原理。

我们将其命名为\textbf{最小空间编码原理 (Minimal Spatial Encoding Principle)}。该原理断言：任何一个在离散、混沌的计算场中传播的\textbf{相干信息模式 (Coherent Information Pattern)}，都会受到两种协同作用的动力学机制的约束，使其自发地、非线性地被局域化在一个有限的时空尺度内。

为了分析这一原理，我们将一个信息模式（Pattern）分为\textbf{整体（the whole）}和\textbf{局部（the part）}，其动力学由以下两个基本假设所支配：


\subsection{机制一：自生场的``内聚效应'' (The Cohesive Effect of the Self-Generated Field)}

\begin{itemize}
    \item \textit{假设1 (干涉 / 时间膨胀效应):} 一个周期性信息的传播，在其同频的、整体信息周围，会因为\textbf{干涉效应}或\textbf{“带宽占用”}而导致其有效空间传播速度变慢。这个作用在物理上，类似于\textbf{克尔效应（Kerr Effect）}的自聚焦。
    
    \item \textbf{推论:} 根据最小作用量原理，这种由模式自身所产生的``慢速介质''效应，会产生一种强大的\textbf{``自聚焦'' (Self-focusing)}倾向。这个机制主要贡献了模式的\textbf{内聚性 (Cohesion)}，使其信息向中心聚集。
\end{itemize}

\subsection{机制二：混沌背景的``边界效应'' (The Boundary Effect of the Chaotic Background)}

\begin{itemize}
    \item \textit{假设2 (随机反射效应):} 信息在离散的、随机的计算场中传播时，其向外弥散的``波前''会与场中无数个随机的``比特缺陷''（由`$b_0$'混沌所体现）发生复杂的相互作用，从而被\textbf{概率性地反射}回中心。这个机制的物理本质，与\textbf{安德森局域化（Anderson Localization）}现象同源。
    
    \item \textbf{推论:} 这个效应为模式提供了一个有效的\textbf{``概率性反射墙''或``边界'' (Probabilistic Boundary)}，确保了其信息不会无限耗散和逃逸。
\end{itemize}


\subsection{综合效应与物理实在的涌现}

这两个机制协同作用：机制一（自聚焦）从\textbf{内部}提供了强大的\textbf{内聚力}，使模式倾向于聚拢；机制二（混沌反射）从\textbf{外部}提供了坚固的\textbf{约束力}，防止模式无限弥散。

一个稳定的物质粒子（孤子），正是信息的扩散倾向和聚拢效应达成完美动态平衡时，所形成的一个非线性的共振结构。

模式传播的\textbf{方向性 (Directionality)}，则源于其创生事件（如原子跃迁）中，由源头的几何不对称性所赋予的一个确定的、量子化的\textbf{``聚焦式初始动量''}。

深刻的是，\textbf{机制一（自聚焦/时间膨胀效应）}与我们理论中的\textbf{引力}是同一个根本机制，引力是作用在所有计算层上的宏观平均体现。\textbf{宏观物质}因其巨大的平均效应，其行为主要由机制一主导，几乎不受机制二（随机反射）的微观效应影响，因此可以呈现出分散的形态。这两种效应的显著与否，是一个\textbf{尺度依赖 (Scale-Dependent)}的问题。

\section{力的分野与物质构成——从梯度力到束缚力}

\subsection{“荷”的统一与分层定义：一个``信息喷泉''模型}

\begin{itemize}
    \item \textbf{荷的统一物理图像 (Unified Physical Picture of Charge):}  
    在一个信息孤子（如电子或夸克）的中心，存在一个持续的、整体的\textbf{跨计算层信息流}（向上或向下）。这个核心信息流，在\textbf{每一个}计算层上，都会在空间维度上\textbf{向外扩展}，形成一个环状的流场，然后在一定的空间距离之外，\textbf{重新汇聚并下沉/上浮}回相邻的层级，形成一个自洽的、闭合的\textbf{``信息涡环'' (Information Torus)} 或 \textbf{``信息喷泉'' (Information Fountain)}。我们必须强调，这个``信息涡环''描述是在2维空间情景下为了概念清晰而进行的简化描述，物理实在是在三维空间中展开的``3维信息喷泉''。  
    
    我们所说的\textbf{``荷''，就是这个``信息喷泉''在不同计算层级上的宏观表现}。
    
    \item \textbf{荷的分层定义 (Hierarchical Definition of Charges):}
    \begin{itemize}
        \item \textbf{电荷 (Electric Charge):} 是这个``信息喷泉''在\textbf{`$b_0-b_1$'}层级上所形成的\textbf{涡环}。这个涡环可以扩展到\textbf{一定的空间距离}，形成一个空间的梯度场。
    
        \item \textbf{强荷 (Strong Charge):} 是这个``信息喷泉''在\textbf{`$b_1-b_2$'}层级上所形成的\textbf{涡环}。由于受到我们下面描述的``强荷寄生性''公理的约束，这个更高层级的涡环\textbf{无法}脱离其`$b_0-b_1$'的能量基座而长距离存在，因此它被\textbf{囚禁}在一个\textbf{极其短的空间区域}内。
    \end{itemize}
    
    \item \textbf{跨层传输协议 (Cross-Layer Transport Protocols):}
    \begin{itemize}
        \item \textbf{单向性约束 (Unidirectionality Constraint):} 在一个稳定的孤子内部，这个``信息喷泉''的整体流动方向必须是一致的（要么是整体向上的``喷泉''，要么是整体向下的``漏斗''）。
        \item \textbf{带宽量子化 (Bandwidth Quantization):} `$b_0-b_1$'通道的最小信息单位为\textbf{1}；`$b_1-b_2$'通道的最小信息单位为\textbf{1/3}。
    \end{itemize}
\end{itemize}

\subsection{电磁相互作用：一种``自主的梯度力''}

\begin{itemize}
    \item \textbf{电子的定义 (Definition of the Electron):} 电子，是只涉及`$b_0-b_1$'层信息流的、简单的稳定信息孤子。
    \item \textbf{电磁力的本质 (Nature of Electromagnetic Force):} 一个带电荷的孤子（如电子），会在其周围的`$b_0-b_1$'计算层上维持一个静态的信息梯度场（电场）。电磁力就是带电孤子与其他电荷产生的梯度场发生的相互作用。单个电荷产生的梯度场只维持在有限的空间距离，宏观的电场是大量电荷的平均效应。
\end{itemize}

\subsection{强相互作用：一种``寄生的束缚力''}

\begin{itemize}
    \item \textbf{夸克的定义 (Definition of the Quark):} 夸克，是一种跨越`$b_0-b_1-b_2$'三层的复合信息孤子，它同时携带\textbf{电荷}和\textbf{强荷}。
    
    \item \textbf{强荷的根本物理学：从``寄生性''到``旋转对称性''}
    \begin{enumerate}
        \item \textbf{第一性公理——强荷寄生性 (Axiom of Parasitic Existence):} 强荷场（`$b_1-b_2$'层的信息流）的存在，必须以一个活跃的`$b_0-b_1$'信息流作为其\textbf{动力学基座}。
        \item \textbf{推论1 (短程性):} 作为该公理的直接后果，强荷场是一个被动力学定律所强制的\textbf{短程场}。
        \item \textbf{推论2 (旋转效应):} 一个无法扩散的场，其唯一稳定的动力学形态，就是将其全部``计算活动''\textbf{内化}为一个局域的、闭合的\textbf{``内部旋转''或``相位''}。
        \item \textbf{推论3 (旋转群对称性):} 描述这种``内部相位旋转''的数学语言，\textbf{必然是一个旋转群}。根据观测到的强子构成规则，我们发现这个群的性质与\textbf{`$U(1)$'旋转群}完美吻合，其代数等价于\textbf{“模1算术”}，在此代数下，任何整数值的净旋转都等价于零 (`$n \pmod{1} = 0$')。
    \end{enumerate}
    
    \item \textbf{强子的构成法则 (The Law of Hadron Formation):}
    \begin{itemize}
        \item \textbf{存在性量子化公理 (Axiom of Quantized Existence):} 任何一个能够作为独立的、自洽的实体稳定存在的``信息模式''（无论是轻子还是强子），其总\textbf{电荷}必须是基本单位的\textbf{整数倍}。
        \item \textbf{荷的守恒约束 (Charge Conservation Constraint):} 任何一个稳定的夸克模式，其\textbf{电荷与强荷之和的绝对值}，必须等于`$b_0-b_1$'通道的基础信息单位\textbf{1}。$|$电荷 + 强荷$| = 1$。
        \item \textbf{夸克的荷分配:} 基于以上公理和约束，第一代只存在四种稳定的夸克模式：
        \begin{enumerate}
            \item \textbf{下夸克 (d):} 电荷`-1/3'，强荷`-2/3'
            \item \textbf{上夸克 (u):} 电荷`+2/3'，强荷`+1/3'
            \item \textbf{反下夸克 ($\bar{d}$):} 电荷`+1/3'，强荷`+2/3'
            \item \textbf{反上夸克 ($\bar{u}$):} 电荷`-2/3'，强荷`-1/3'
        \end{enumerate}
        \item \textbf{强荷中性的涌现 (Emergence of Strong-Charge Neutrality):} 作为一个必然的\textbf{定理}，任何一个满足``存在性量子化公理''（即总电荷为整数）的强子，其内部所有夸克的\textbf{总强荷}在`$U(1)$'的意义下，\textbf{必然为零}。（例如，质子`$uud$'的总强荷为 `$(+1/3) + (+1/3) + (-2/3) = 0 \pmod{1}$'）。
    \end{itemize}
    
    \item \textbf{核力的本质 (Nature of the Nuclear Force):}
    \begin{itemize}
        \item 一个孤立的核子是\textbf{强荷中性}的。
        \item 但当两个核子极其接近时，它们各自内部的\textbf{信息回路发生重叠}。这导致原本在核子内部完美循环的``强荷''信息流发生\textbf{``泄漏''}，在两个核子之间形成一个全新的、\textbf{短程的、跨核子的信息流回路}。这个新形成的回路所产生的强大吸引力，就是``核力''的根本来源。这个机制与分子间的\textbf{范德华力}在动力学上同构。
    \end{itemize}
    
    \item \textbf{夸克禁闭 (Quark Confinement):}
    \begin{itemize}
        \item \textbf{内在原因：} 禁闭的根本原因，是任何一个稳定的强子（重子或介子），其内部所有的夸克，都必须\textbf{``寄生''在同一个`$b_0-b_1$'层的基础信息流结构之上}。一个孤立的夸克，其`$b_0-b_1$'基础结构是不完备的，违反了存在性量子化公理 ，因此无法独立稳定存在。
        \item \textbf{动力学表现：} 夸克之间的``强荷力''极其巨大。任何试图将一个夸克从一个强荷中性的核子中分离出来的尝试，都需要注入巨大的能量。由于夸克的尺度极小，任何能够作用于它的探针，都不可避免地会向系统注入\textbf{远超创生阈值的能量}，其最终结果不是``分离''，而是从真空中创生出新的夸克对，形成新的、满足所有守恒律和中性要求的稳定强子。
    \end{itemize}
\end{itemize}

\section{弱相互作用——孤子的``衰变''与``嬗变''}

在《计算实在论》的框架下，弱相互作用的本体论地位被根本性地重构。它不再被视为一种与其他三种力并列的``相互作用力''，而被理解为``信息孤子''自身的一种\textbf{内在的、动力学不稳定性}。

\subsection{弱相互作用的本质：一个``拓扑重组''的过程}

\begin{itemize}
    \item \textbf{定义 (Definition):} 弱相互作用，是描述一个亚稳态``信息孤子''（如中子），因其\textbf{内因（内在的结构不完美）}与\textbf{外因（与外部环境的持续能量/信息交换，如`$b_0$'混沌涨落或高能碰撞）}共同作用下，被触发并经历一次\textbf{整体``拓扑重组'' (Holistic Topological Reconfiguration)}的动力学过程。
    \item \textbf{核心机制 (Core Mechanism):} 这是一个内部bit信息从有序的孤子结构中\textbf{``溃散''并``破裂''}的过程，其最终目的是让系统演化到总能量更低、结构更稳定的新孤子组合。它本质上是一个\textbf{过程 (Process)}，而非一个瞬时事件。
\end{itemize}

\subsection{“弱力”的再诠释：内部信息的``扩散驱动力''}

\begin{itemize}
    \item \textbf{定义 (Definition):} 在我们的模型中，``弱力''并非一种传统意义上的、描述孤子间相互作用的``力''。它更应被理解为，在那场``拓扑重组''过程中，被孤子结构所``禁闭''的内部信息（能量与动量）向外\textbf{扩散的驱动力}。
    \item \textbf{物理图像 (Physical Picture):} 它就像一个高压容器（亚稳态孤子）突然破裂，其内部的高压气体（被禁闭的信息）向外猛烈喷发。这股``喷发''的驱动力，就是``弱力''的动力学本质。
\end{itemize}

\subsection{W/Z玻色子的本质：高能的``中间共振态''}

\begin{itemize}
    \item \textbf{定义 (Definition):} W/Z玻色子，是在那场剧烈的``拓扑重组''事件中，由被暂时``解禁''的巨大内部能量所形成的、\textbf{真实的、但转瞬即逝的``高能中间共振态'' (High-Energy Intermediate Resonance)}。
    \item \textbf{质量的再解释 (Reinterpretation of Mass):} 像`$80.4~\text{GeV}$'这样的数值，可以被理解为一种\textbf{``瞬时的有效静止质量''}。它代表了一个破溃过程中的碎片，一个\textbf{``临时信息包''}，在其彻底瓦解之前，所能\textbf{稳定束缚的能量上限}。任何超过这个能量上限的注入，都会以动能形式存在或者导致其分裂成独立的``信息包''。
\end{itemize}

\subsection{弱相互作用的核心``签名''}

\begin{itemize}
    \item \textbf{味道的改变与CKM矩阵 (Flavor Changing \& CKM Matrix):}
    \begin{itemize}
        \item \textbf{现象:} 弱相互作用是唯一能改变夸克和轻子``味道''（代际）的力，且这种跨代转变的概率由一个复杂的CKM矩阵所描述。
        \item \textbf{理论解释:}
        \begin{itemize}
            \item \textbf{``味道''的本质:} 不同的``味道''（如u vs s夸克），对应了具有\textbf{不同``跨层信息流''拓扑结构}的孤子。
            \item \textbf{改变``味道'':} ``弱相互作用''作为一场剧烈的``拓扑重组''，其过程自然包含了重新``布线''这些跨层信息流的可能性，因此它是唯一能改变``味道''的机制。
            \item \textbf{CKM矩阵的起源:} CKM矩阵的数值是在``拓扑重组''这个复杂的动力学过程中，从一个初始拓扑构型（如s夸克），跃迁到另一个可能的末态拓扑构型（如u夸克）的\textbf{``跃迁概率幅''}。这些数值，反映了不同拓扑结构之间的\textbf{``可转换性''或``相似度''}。
        \end{itemize}
    \end{itemize}
\end{itemize}

\section{一个统一的猜想——中微子、宇称不守恒与暗物质的共同起源}

在建立了我们关于物质局域化和基本相互作用的动力学模型之后，我们发现，这个框架的内在逻辑，以一种近乎不可避免的方式，将我们引向了对宇宙三个最深刻谜题——中微子的本质、宇称不守恒的起源、以及暗物质的存在——的一个统一的、虽然尚属猜想，但逻辑上高度自洽的解释。

本章旨在分享这个探索性的思路。我们并不声称这是最终的答案，而是希望展现，当我们严格地遵循``分层计算''这个核心公理时，一幅出人意料的宇宙图景，可能会在我们面前徐徐展开。

\subsection{中微子之谜：一个``平行家族''的信使？}

\noindent\textbf{观察到的谜题：}\\
中微子展现出一系列极其独特的``反常''特性：1) 它近乎``幽灵''，只参与弱相互作用和引力，完全不参与电磁相互作用（电中性）；2) 它的质量极其微小，远低于所有其他物质粒子。为什么存在这样一种几乎与我们可见世界``绝缘''的粒子？

\noindent\textbf{一个可能的解答（我们的猜想）：}\\
我们猜想，中微子之所以如此独特，是因为它根本不属于我们所熟知的``可见物质家族''。

\begin{itemize}
    \item \textbf{`$b_0$'层起振的``可见家族''：} 我们提出，我们宇宙中所有可被电磁力感知的物质（电子、夸克），其共同特征在于，它们的存在都\textbf{``构筑在`$b_0$'层起振''}。它们的稳定结构，都深刻地依赖于与最底层`$b_0$'混沌的直接交互，这个交互的宏观体现就是``电荷''。

    \item \textbf{中微子——`$b_1$'层起振的``微光家族''：} 与此相对，我们猜想，中微子家族是第一个、也是最简单的、一个纯粹的\textbf{```$b_1$'层起振''的粒子家族}。这个猜想，为中微子的核心谜题提供了一个统一的动力学解释：
    \begin{enumerate}
        \item \textbf{它为何电中性？} \\
        因为中微子的核心稳定结构，\textbf{没有一个`$b_0-b_1$'的跨层信息流循环}。它的`$b_1-b_2$'循环对`$b_0-b_1$'层造成的影响极其微弱，这可能是因为两者\textbf{共振频率差别太大}；同时，在宏观上，宇宙背景中的\textbf{中微子和反中微子密度相等}，其微弱的次级影响在统计上会相互抵消。

        \item \textbf{它为何质量极小？} \\
		粒子的静止质量，主要源于其跨层构建时，为了克服不同计算层\textbf{``计算核''的尺寸``失谐''}而必须付出的结构成本。中微子，作为一个独立的$b_1-b_2$层的孤子，相对于一个$b_0-b_1$层孤子，其``失谐''成本更小，所以只需一个更小的耦合结构就能稳定存在，因此其静止质量小。
    \end{enumerate}
\end{itemize}

\subsection{宇称不守恒之谜：一个动态的、局域的相对性效应}

\noindent\textbf{观察到的谜题：}\\
弱相互作用——这个与中微子紧密相关的力——最大程度地破坏了``镜像对称性''（宇称）。实验证实，宇宙在弱相互作用的层面上，是一个坚定的``左撇子''。这种深刻的、根本性的不对称性，其起源是什么？

\noindent\textbf{一个可能的解答（我们的猜想）：}\\
我们猜想，宇称不守恒并非一个全局性的、固定的``宇宙规则''，而是一个\textbf{动态的、局域的、源于我们理论中``分层计算''这个核心公理的必然涌现}。其本质，是一种深刻的\textbf{``层级相对性原理'' (Principle of Hierarchical Relativity)}。

\begin{itemize}
    \item \textbf{层级相对性原理：} 我们假设，`$b_0$'层和`$b_1$'层是两个\textbf{不同的、平等的动力学基准}。它们之间，在每一个空间区域内，都存在着一种\textbf{时刻演化的、内禀的``相对旋转''}。手性（Chirality），就是一个粒子，与`$b_0$'层和`$b_1$'层旋转的接近程度。

    \item \textbf{对称与不对称的涌现 (The Emergence of Symmetry and Asymmetry):} \\
    一个粒子是否表现出手性对称，完全取决于其自身的\textbf{``计算层级构成''}及其与这两个不同动力学基准的耦合关系。
    \begin{enumerate}
        \item \textbf{可见物质（`$b_0$'家族）——手性对称的``混合态''：}\\
        一个``正常''的、可见的粒子（如电子），作为一个\textbf{```$b_0$'层起振''}的实体，其稳定存在是一个横跨`$b_0$'和`$b_1$'的\textbf{``混合态''}。因此，它必然存在两种最基本的动力学模式来定义其自旋/手性：一种是\textbf{``更接近$b_0$''}的模式（我们定义为右手性），另一种是\textbf{``更接近$b_1$''}的模式（我们定义为左手性）。拥有这两种稳定状态，使得可见粒子整体上表现为\textbf{手性对称的}。

        \item \textbf{中微子（`$b_1$'家族）——手性不对称的``纯粹态''：}\\
        中微子，作为一个纯粹的\textbf{```$b_1$'层起振''}的实体，它\textbf{没有``更接近`$b_0$'''这种状态可以选择}。在`$b_0$'家族的粒子看来，它唯一的、能够稳定存在的动力学模式，就是完全与它所在的`$b_1$'基准的``相对旋转''特性保持\textbf{同步}。因此，从我们`$b_0$'观察者的视角看，\textbf{中微子，始终是与那个``相对旋转''的`$b_1$'基准是绑定的}，并因此只表现出一种手性（即我们定义的``左手性''）。
    \end{enumerate}

    \item \textbf{手性场的动态与局域性 (The Dynamic and Local Nature of the Chiral Field):}\\
    我们进一步推断，`$b_0$'与`$b_1$'之间的``相对旋转''，是一个在时空中\textbf{动态变化}的局域场。每一个手性粒子（如中微子）的创生和衰变，都会通过一种``作用力与反作用力''的机制，向其周围的`$b_0-b_1$'背景场施加一个``反向手性''的贡献。它是一种\textbf{宏观的、时空局域的、动态的平衡态}。而这个场，很可能跟我们之前论述的纠缠的背景场是同一个实体。
\end{itemize}

\subsection{暗物质的必然存在性}
\noindent\textbf{从上述猜想引出的必然推论：}\\
一旦我们接受了``粒子家族可以由不同计算层起振''这个想法，并用它成功地解释了中微子的存在和弱力的不对称性，那么一个更加宏大和不可避免的推论就会自然而然地出现：

\begin{itemize}
    \item 如果存在`$b_0$'层起振的``可见世界''，和`$b_1$'层起振的``微光世界''（中微子），那么，\textbf{必然存在在更高的计算层起振的、更``黑暗''的、更``隐形''的粒子家族。}

    \item \textbf{暗物质的终极身份：} 我们因此提出终极猜想——我们宇宙中那占据了绝大部分物质质量的\textbf{暗物质}，其主体，正是这些由`$b_2$', `$b_3$'或更高计算层独立``起振''的、与我们可见世界几乎完全``通信协议不兼容''的\textbf{``暗黑粒子家族''}。它们不参与电磁力、强力、甚至可能连弱力都不参与。它们与我们共享的、唯一的、不可避免的连接，就是\textbf{引力}。
\end{itemize}

\subsection{一个更大胆的猜想——中微子的引力}
我们猜测，中微子的惯性质量和引力质量比，与`$b_0$'起点的物质可能存在差别。

我们所定义的惯性质量和引力质量，其标尺，引力常数，是在\textbf{“`$b_0$'层为共同共振起点”}的基础上建立的。中微子这种`$b_1$'层起振的粒子，其惯性质量和引力质量的比例，可能出现了与正常物质相比的\textbf{``显著分离''}。因为`$b_1$'层更大的计算核对下层的作用，可能导致中微子按标准定义下的引力质量大于惯性质量，或者说它有更大的引力常数。

作为类比，有一个众周知的、违反这个惯性质量和引力质量比的实体，是光子。

\section*{总结}
至此，我们已经完成了一次基于《计算实在论》的探索性思考。我们试图从第一性原理出发，为粒子物理学的核心谜团提供一个统一的、内在自洽的解释。如果这个思考过程本身，能为您带来任何有益的启发，我们便心满意足。


\end{document}